\chapter{Unemployment: Stocks and Flows}\label{ch:b1:unemployment-flows}

In an average month of 2018, American employers hired 5.71 million people
and 5.51 million left or lost a job. % figdata: jolts_stats
The annual unemployment rate fell from $4.36\%$ in 2017 to $3.9\%$ in 2018.
% figdata: us_unemp_points
The rate is a stock; behind it are flows many times its monthly change.
This chapter models unemployment as the balance of those flows, asks why
wages do not fall until every job seeker is hired, and reads the evidence
on how benefits change the flows.

\begin{openingchart}{Hires and total separations, United States, monthly,
2001--2025 (millions, seasonally adjusted, total nonfarm).}{U.S. Bureau of Labor
Statistics, Job Openings and Labor Turnover Survey}
\begin{tikzpicture}
\begin{axis}[omtime, xlabel={Year}, ylabel={Millions a month}, xmin=2001, xmax=2026, ymin=0, ymax=12,
  width=0.72\linewidth, height=0.4\linewidth, legend style={font=\footnotesize, at={(0.5,-0.24)}, anchor=north,
  legend columns=2}]
  \addplot+[thick, no marks] table[x=t, y=HI, col sep=comma] {figdata/bachelor-1/22-unemployment-flows/out/jolts_flows.csv};
  \addplot+[thick, no marks] table[x=t, y=TS, col sep=comma] {figdata/bachelor-1/22-unemployment-flows/out/jolts_flows.csv};
  \legend{Hires, Separations}
\end{axis}
\end{tikzpicture}
\end{openingchart}

The two lines nearly coincide; the gap between them is the monthly change
in employment. In 2009 separations exceeded hires by about 0.4 million a
month, 4.33 million against 3.9; in April 2020 11.53 million jobs ended,
9.17 million of them layoffs, while 4.03 million people were hired.
% figdata: jolts_stats
% figdata: jolts_2020
Recessions show up as a fall in hiring as much as a rise in separations.

\begin{recall}
The unemployment and participation rates and the survey that measures
them were met at school, as were the flows into and out of unemployment,
the natural rate and the reasons wages may not clear the market. \omterm{def:b1:factor-markets:monopsony}{Monopsony}
and wage floors are in \cref{ch:b1:factor-markets}.
\end{recall}

\section{Flows and the steady state}\label{sec:b1:unemployment-flows:flows}

Take a labour force of fixed size, normalised to one, with a share $u$
unemployed. Each month a share $s$ of the employed lose or leave their job
and become unemployed, and a share $f$ of the unemployed find a job.

\begin{definition}[Separation rate, job-finding rate]\label{def:b1:unemployment-flows:rates}
\index{separation rate}\index{job-finding rate}
The \emph{separation rate} $s$ is the share of the employed who become
unemployed in a period; the \emph{job-finding rate} $f$ is the share of the
unemployed who find a job in a period.
\end{definition}

The unemployment rate then changes as
\[
\dot u = s(1 - u) - fu .
\]

\begin{definition}[Steady-state unemployment rate]\label{def:b1:unemployment-flows:steady-state-rate}
\index{steady-state unemployment rate}
The \emph{steady-state unemployment rate} is the rate at which inflows equal
outflows, $s(1 - u^*) = fu^*$:
\[
u^* = \frac{s}{s + f}.
\]
\end{definition}

\begin{proposition}[Convergence]\label{prop:b1:unemployment-flows:convergence}
With constant $s$ and $f$, $u(t) = u^* + (u(0) - u^*)e^{-(s + f)t}$: the gap closes at the rate $s + f$.
\end{proposition}

\begin{proof}
$\dot u = s - (s + f)u = -(s + f)(u - u^*)$, a linear equation whose solution is the one stated.
\end{proof}

\begin{omfigure}
\begin{tikzpicture}
\begin{axis}[omrecord, grid=none, xlabel={Unemployment rate $u$}, ylabel={Flow (share of labour force a month)},
  xmin=0, xmax=0.15, ymin=0, ymax=0.045, width=0.56\linewidth, height=0.4\linewidth, scaled ticks=false,
  xticklabel style={/pgf/number format/fixed}, yticklabel style={/pgf/number format/fixed}]
  \addplot[supply, domain=0:0.15] {0.01917*(1 - x)};
  \addplot[demand, domain=0:0.15] {0.30*x};
  \addplot[guide] coordinates {(0.06006,0) (0.06006,0.01802)};
  \node[font=\footnotesize, anchor=south] at (axis cs:0.128,0.0168) {$s(1 - u)$};
  \node[font=\footnotesize, anchor=west] at (axis cs:0.1,0.035) {$fu$};
  \node[font=\footnotesize, anchor=north west] at (axis cs:0.061,0.004) {$u^* = 6.0\%$};
\end{axis}
\end{tikzpicture}

{\small The \omterm{def:b1:solow-model:steady-state}{steady state} of the weekend problem: the inflow into
unemployment $s(1 - u)$ (red) and the outflow $fu$ (blue), with $s = 0.01917$ and
$f = 0.30$ a month, cross at $u^* = 6.0\%$. Above it the outflow is larger and $u$
falls; below it $u$ rises.}
\end{omfigure}

Monthly rates are large: with $f = 0.30$ and $s$ near $0.02$, $s + f$ is about $0.32$,
and half of any gap closes in a little over two months. The actual rate is
therefore never far from the steady-state rate of the moment; when it moves
for years, it is because $s$ and $f$ have moved. The two rates are not the
natural rate of unemployment met at school, the rate around which
unemployment fluctuates and at which inflation is stable
(\cref{ch:b1:aggregate-supply}); the steady-state rate at the separation
and finding rates of normal, non-recession years is one way to estimate it.

\begin{measure}[The flow data]
Flows of workers between employment, unemployment and inactivity are
measured by matching people in the monthly household survey from one
month to the next: about three-quarters of the households interviewed in a
month were also interviewed the month before. The United States publishes
nine core flows, in a series starting in February 1990. % ledger: b1:bls-flows-method
The hires and separations of the opening chart come from a different
survey, of employers, and include moves from one job to another.
\end{measure}

Which flow moves unemployment in recessions? Since 1948, changes in the
job-finding probability account for three-quarters of the fluctuations of
the American unemployment rate and changes in the probability of leaving
employment for one-quarter; over the last two decades of that study, the
second was negligible (Shimer 2012). % ledger: b1:shimer2012
Unemployment rises in a recession less because more people lose their jobs
than because those without one find a new one more slowly.

Over a longer span the American rate moved far more than in any recent
recession: $24.9\%$ in 1933, $1.2\%$ in 1944, $9.7\%$ in 1982. % figdata: us_unemp_points

\begin{omchart}{Unemployment rate, United States, 1890--2018.}{Jord\`a--Schularick--Taylor
Macrohistory Database, release 6}
\begin{tikzpicture}
\begin{axis}[omtime, xlabel={Year}, ylabel={Per cent of the labour force}, xmin=1890, xmax=2018, ymin=0, ymax=27,
  width=0.72\linewidth, height=0.36\linewidth]
  \addplot+[thick, no marks] table[x=year, y=unemp, col sep=comma] {figdata/bachelor-1/22-unemployment-flows/out/us_unemp.csv};
\end{axis}
\end{tikzpicture}
\end{omchart}

\begin{model}{The flow model of unemployment}
\assumes Constant separation and \omterm{def:b1:unemployment-flows:rates}{job-finding rates}; a fixed labour force;
no flows into or out of inactivity.
\predicts A steady-state rate $u^* = s/(s + f)$, and convergence to it at the rate
$s + f$, fast at monthly frequency.
\tested In American data, most of the cyclical movement comes from the
\omterm{def:b1:unemployment-flows:rates}{job-finding rate} (Shimer 2012). % ledger: b1:shimer2012
\breaks Flows to and from inactivity are large and ignored; rates vary with
the duration of unemployment and across workers, so the averages hide
long spells.
\end{model}

\section{Duration}\label{sec:b1:unemployment-flows:duration}

\begin{definition}[Expected unemployment duration]\label{def:b1:unemployment-flows:duration}
\index{unemployment duration}
With a constant monthly \omterm{def:b1:unemployment-flows:rates}{job-finding rate} $f$, the number of months until a
job is found is geometric, and the \emph{expected duration} of a spell is $1/f$.
\end{definition}

With $f = 0.30$ the expected spell is 3.33 months. But $f$ is not the same for
everyone: if it falls the longer a person has been unemployed --- because
skills or contacts fade, or because employers read a long spell as a bad
signal --- there is \emph{duration dependence}, and a minority of long spells
makes up much of total unemployment. With two groups, one finding jobs at
$0.4$ a month and one at $0.05$, each with half of the spells, the average
spell lasts $(2.5 + 20)/2 = 11.25$ months, and at any moment the long-term group is
$20/22.5$, about $89\%$, of the unemployed.
% computed: code/bachelor-1/22-unemployment-flows/exercises/solutions.py

\section{Why wages stay above the clearing level}\label{sec:b1:unemployment-flows:wages}

If wages fell until everyone who wants a job had one, unemployment would
only be the time spent between jobs. Three mechanisms keep the wage above
that level.

\subsection*{Efficiency wages}

Suppose effort cannot be enforced directly. A worker who works hard bears
a cost $e$ (per month, in money terms); a worker who shirks is caught with
probability $q$ per month and fired. Unemployed, a worker receives $z$
(benefits and home production) and finds a job at the rate $f$; the
interest rate is $r$.

\begin{definition}[Efficiency wage, no-shirking condition]\label{def:b1:unemployment-flows:efficiency-wage}
\index{efficiency wage}\index{no-shirking condition}
An \emph{efficiency wage} is a wage above the market-clearing level, paid
because it raises productivity --- through effort, lower turnover or
better recruits. The \emph{no-shirking condition} is the lowest wage at which
a worker prefers effort to shirking.
\end{definition}

\begin{proposition}[The no-shirking wage]\label{prop:b1:unemployment-flows:nsc}
In continuous time, with values $V_E$ (employed, working), $V_S$ (employed,
shirking) and $V_U$ (unemployed), workers do not shirk if and only if
\[
w \ge z + e + \frac{e}{q}(r + s + f).
\]
\end{proposition}

\begin{proof}
The asset equations are $rV_E = w - e + s(V_U - V_E)$, $rV_S = w + (s + q)(V_U - V_S)$ and
$rV_U = z + f(V_E - V_U)$. Workers exert effort if $V_E \ge V_S$; the firm pays the
lowest such wage, so $V_E = V_S$. Subtracting the first two equations with
$V_E = V_S$ gives $e = q(V_E - V_U)$: the rent from the job must be $e/q$. Subtracting the
third from the first, $r(V_E - V_U) = w - e - z - (s + f)(V_E - V_U)$, so
$w = z + e + (r + s + f)\,e/q$.
\end{proof}

The wage must exceed the worker's alternative by a rent, and the rent is
larger when detection is rare (small $q$) and when a fired worker finds a
new job quickly (large $f$). In a \omterm{def:b1:solow-model:steady-state}{steady state} $f = s(1 - u)/u$, so the condition
becomes $w = z + e + (e/q)(r + s/u)$: the lower unemployment, the higher the wage
needed. Unemployment is what makes the threat of dismissal bite.

\begin{omfigure}
\begin{tikzpicture}
\begin{axis}[omrecord, grid=none, xlabel={Employment (share of labour force)}, ylabel={Wage},
  xmin=0.8, xmax=1, ymin=0, ymax=12, width=0.56\linewidth, height=0.4\linewidth, samples=80,
  xtick={0.8,0.85,0.9,0.95,1}, xticklabel style={/pgf/number format/fixed}]
  \addplot[supply, domain=0.8:0.985] {1.2 + 0.5 + 5*(0.005 + 0.01917/(1 - x))};
  \addplot[demand, domain=0.8:1] {3.3209 + 30*(0.94 - x)};
  \addplot[guide] coordinates {(0.94,0) (0.94,3.3209)};
  \addplot[guide] coordinates {(1,0) (1,12)};
  \node[font=\footnotesize, anchor=south east] at (axis cs:0.98,9) {NSC};
  \node[font=\footnotesize, anchor=south west] at (axis cs:0.81,8.2) {labour demand};
\end{axis}
\end{tikzpicture}

{\small The \omterm{def:b1:unemployment-flows:efficiency-wage}{no-shirking condition} of the weekend problem (red) rises steeply
as employment approaches the whole labour force; with labour demand
(blue) it sets employment at $94\%$ and the wage at $3.32$ thousand dalens a
month. Unemployment is involuntary: the unemployed would work for less,
but a lower wage would not be credible.}
\end{omfigure}

\begin{record}[Five dollars a day]
In 1914 Ford introduced the five-dollar day, roughly doubling the pay of
its production workers. The decision followed labour problems of the kind
efficiency-wage models stress --- high turnover and absenteeism --- and it
was followed by long queues for Ford jobs and significant increases in
productivity and profits (Raff and Summers 1987). % ledger: b1:raff1987
A firm paid more than it needed to hire, and gained by it.
\end{record}

\begin{model}{Efficiency wages}
\assumes Effort that is costly to the worker and imperfectly observed;
dismissal as the penalty; workers who value the job by its rent over
unemployment.
\predicts Wages above the clearing level; involuntary unemployment; lower
turnover and more effort at high-wage firms.
\tested Ford's wage of 1914 brought queues of applicants, lower turnover
and higher productivity (Raff and Summers 1987). % ledger: b1:raff1987
\breaks It does not explain why wages fall so little in deep recessions,
when the rent from a job rises on its own as $f$ falls.
\end{model}

\subsection*{Bargaining and minimum wages}

Where unions or individual workers bargain, the wage splits the surplus of
the match between the firm and the worker; a stronger bargaining position
raises the wage above the worker's alternative and lowers the number of
jobs firms open. A legal minimum above the clearing wage has the effects of
\cref{ch:b1:factor-markets}: in a competitive market it lowers employment,
in a \omterm{def:b1:factor-markets:monopsony}{monopsony} it can raise it.

\section{The Beveridge curve}\label{sec:b1:unemployment-flows:beveridge}

\begin{definition}[Beveridge curve]\label{def:b1:unemployment-flows:beveridge}
\index{Beveridge curve}
The \emph{Beveridge curve} is the downward-sloping relation between the
vacancy rate and the unemployment rate: when firms look for many workers,
few workers look for jobs.
\end{definition}

Along the curve, an economy moves with the cycle: in a boom vacancies rise
and unemployment falls. A shift of the curve outwards means more vacancies
\emph{and} more unemployment at once: jobs and workers are matched less
well, because skills, places or sectors do not fit. The matching function
behind the curve is a subject of the second university year.

\begin{omchart}{The Beveridge curve of the European Union: job vacancy rate
against unemployment rate, quarterly, 2010--2025 (seasonally adjusted).}{Eurostat
(jvs\_q\_nace2, une\_rt\_q)}
\begin{tikzpicture}
\begin{axis}[omscatter, xlabel={Unemployment rate (\%)}, ylabel={Job vacancy rate (\%)}, xmin=5.5, xmax=12,
  ymin=0.8, ymax=3.4, width=0.6\linewidth, height=0.42\linewidth]
  \addplot[omDef, thin, sharp plot, mark=*, mark size=1.2pt] table[x=u, y=vac, col sep=comma] {figdata/bachelor-1/22-unemployment-flows/out/beveridge.csv};
  \addplot[only marks, mark=none, nodes near coords, point meta=explicit symbolic, nodes near coords style={font=\scriptsize, anchor=south west}]
    table[x=u, y=vac, meta=time, col sep=comma] {figdata/bachelor-1/22-unemployment-flows/out/beveridge_labels.csv};
\end{axis}
\end{tikzpicture}
\end{omchart}

From early 2013 to the end of 2019 the Union moved up along its curve, to
the left,
from $11.7\%$ unemployment and a vacancy rate of $1.2\%$ to $6.7\%$ and $2.3\%$.
In 2022 the vacancy rate reached $3.1\%$ with unemployment at $6.2\%$; by the end
of 2025 it was back to $2.1\%$ with unemployment at $6\%$. % figdata: beveridge_points
Whether the curve shifted after 2020, or the economy only travelled along
it fast, is the kind of question the matching function is built to answer.

\section{Benefits and hiring rules as evidence}\label{sec:b1:unemployment-flows:benefits}

Unemployment benefits raise $z$: in the flow model they lower the \omterm{def:b1:unemployment-flows:rates}{job-finding
rate} if they make the unemployed search less or accept fewer offers. Rules
that make dismissal costly lower the \omterm{def:b1:unemployment-flows:rates}{separation rate}, but also make firms
hire more cautiously, lowering $f$; with both falling in proportion, $u^*$ does
not change, but spells are longer and adjustment slower.

\begin{debate}[Do longer benefits lengthen unemployment?]
In Germany the maximum duration of benefits depends on exact age, so
workers just above and just below an age threshold can be compared. Over
20 years of such comparisons, the effect of an extra month of benefits on
time out of work is at best somewhat smaller in recessions, while the share
of workers who exhaust their benefits rises sharply; the cost in
incentives of longer benefits is lower in recessions (Schmieder, von
Wachter and Bender 2012). % ledger: b1:schmieder2012
In Austria, extending benefits from 20 to 30 weeks lowered \omterm{def:b1:unemployment-flows:rates}{job-finding rates}
in the first 20 weeks by $5$ to $9\%$, but a severance payment of two months'
earnings, which changes no incentive at the margin, lowered them by $8$ to
$12\%$: much of the effect of benefits may come from cash in hand, not from
distorted incentives (Card, Chetty and Weber 2007). % ledger: b1:card2007
Each design has a weakness: age thresholds compare workers near one age
only, and severance pay goes to workers whose jobs and firms differ in
other ways. The measured effects are clear; their reading is not.
\end{debate}

\section*{Exercises}\label{sec:b1:unemployment-flows:exercises}

\begin{exercise}[$\star$]\label{exo:b1:unemployment-flows:ustar}
With $s = 2\%$ and $f = 28\%$ a month, compute the \omterm{def:b1:unemployment-flows:steady-state-rate}{steady-state unemployment rate}.
\end{exercise}

\begin{exercise}[$\star$]\label{exo:b1:unemployment-flows:dur}
With $f = 25\%$ a month, what is the expected duration of a spell?
\end{exercise}

\begin{exercise}[$\star$]\label{exo:b1:unemployment-flows:read}
The American unemployment rate was $24.9\%$ in 1933, $1.2\%$ in 1944, $9.7\%$ in
1982 and $9.25\%$ in 2009. % figdata: us_unemp_points
Using the chart of the American rate in \cref{sec:b1:unemployment-flows:flows},
say which flow, $s$ or $f$, you would expect to have moved most in 1933 and in
2009, and why the 1944 rate was so low.
\end{exercise}

\begin{exercise}[$\star$]\label{exo:b1:unemployment-flows:predict}
In the economy of \cref{exo:b1:unemployment-flows:ustar}, the \omterm{def:b1:unemployment-flows:rates}{job-finding rate}
falls to $20\%$. Predict the new steady-state rate.
\end{exercise}

\begin{exercise}[$\star\star$]\label{exo:b1:unemployment-flows:half}
With $s + f = 0.30$ a month, how long does it take to close half of a gap
between $u$ and $u^*$?
\end{exercise}

\begin{exercise}[$\star\star$]\label{exo:b1:unemployment-flows:which}
In month A, $s = 0.015$ and $f = 0.30$; in month B, $s = 0.017$ and $f = 0.22$. Compute $u^*$ in
each, and the $u^*$ that each change alone would give. Which flow moved
unemployment more?
\end{exercise}

\begin{exercise}[$\star\star$]\label{exo:b1:unemployment-flows:nsw}
With $z = 1$, $e = 0.4$, $q = 0.2$, $r = 0.004$, $s = 0.02$ and $f = 0.28$ (monthly), find the
no-shirking wage.
\end{exercise}

\begin{exercise}[$\star\star$]\label{exo:b1:unemployment-flows:firing}
A firing cost lowers both $s$ and $f$ by $20\%$ in the economy of
\cref{exo:b1:unemployment-flows:ustar}. Predict $u^*$ and the \omterm{def:b1:solow-model:speed}{half-life} of
convergence.
\end{exercise}

\begin{exercise}[$\star\star$]\label{exo:b1:unemployment-flows:youth}
In 2023 the unemployment rate of 15--24-year-olds and of all 15--74-year-olds
was $5.9\%$ and $3.1\%$ in Germany, $28.7\%$ and $12.2\%$ in Spain, $22.1\%$ and $7.7\%$
in Sweden, and $11.4\%$ and $2.8\%$ in Poland. % figdata: youth
Compute the ratios. Why are young people's rates higher, in flow terms?
\end{exercise}

\begin{exercise}[$\star\star\star$]\label{exo:b1:unemployment-flows:solve}
Solve $\dot u = s(1 - u) - fu$ for constant $s$ and $f$, and show that $u$ converges to $u^*$
at the rate $s + f$ from any $u(0)$ in $[0, 1]$.
\end{exercise}

\begin{exercise}[$\star\star\star$]\label{exo:b1:unemployment-flows:three}
Add inactivity: the unemployed leave the labour force at the rate $x$, and
the inactive enter unemployment at the rate $y$, with no other flows to or
from inactivity. Derive the steady-state shares and the unemployment rate.
\end{exercise}

\begin{exercise}[$\star\star\star$]\label{exo:b1:unemployment-flows:refute}
What observation would refute the claim that recessions raise
unemployment mainly through more separations?
\end{exercise}


\section{Problem: Brenland after a Factory Closes}\label{sec:b1:unemployment-flows:problem}

\begin{problem}[{Weekend problem --- a layoff and a reform}]\label{pb:b1:unemployment-flows:brenland}
In 2025 Brenland (invented) had a labour force of 4,645 thousand and
employment of 4,366 thousand. Rates are monthly. Take a \omterm{def:b1:unemployment-flows:rates}{job-finding rate}
$f = 0.30$ and calibrate the \omterm{def:b1:unemployment-flows:rates}{separation rate} so that 2025 is a \omterm{def:b1:solow-model:steady-state}{steady state}.

\medskip\noindent\textbf{Part I --- Flows and the steady state.}
\begin{enumerate}
  \item Compute the number unemployed and the unemployment rate.
  \item Calibrate $s$.
  \item Compute the monthly inflow into and outflow from unemployment.
  \item Compute the expected duration of a spell.
  \item Compute the \omterm{def:b1:solow-model:speed}{speed of convergence} and the \omterm{def:b1:solow-model:speed}{half-life}.
\end{enumerate}

\medskip\noindent\textbf{Part II --- A factory closes.} In one month $2\%$ of the
employed lose their jobs at once; $s$ and $f$ are unchanged.
\begin{enumerate}[resume]
  \item How many lose their jobs?
  \item Compute the unemployment rate just after.
  \item Write the path of $u$.
  \item Compute $u$ after 1, 3, 6 and 12 months.
  \item Why does the rate return so fast?
  \item What would make the return slower?
\end{enumerate}

\medskip\noindent\textbf{Part III --- Efficiency wages at the new employer.} A
firm that hires the laid-off workers cannot observe effort. Take $z = 1.2$,
$e = 0.5$ (thousand dalens a month), $q = 0.1$ and $r = 0.005$.
\begin{enumerate}[resume]
  \item Compute the no-shirking wage.
  \item What rent does it leave the worker over $z + e$?
  \item How does the wage change if $z$ rises to 1.5?
  \item How does it change if $f$ falls to 0.25?
  \item Why does the firm not cut the wage when the unemployed queue for its jobs?
  \item What would a better monitoring technology (higher $q$) do?
\end{enumerate}

\medskip\noindent\textbf{Part IV --- A benefit reform.} A longer benefit duration
lowers the \omterm{def:b1:unemployment-flows:rates}{job-finding rate} to $0.25$; $s$ is unchanged.
\begin{enumerate}[resume]
  \item Compute the new steady-state rate.
  \item Compute the number unemployed in the new \omterm{def:b1:solow-model:steady-state}{steady state}.
  \item Compute the new expected duration.
  \item Compute the new \omterm{def:b1:solow-model:speed}{half-life}.
  \item By the Austrian evidence, part of the fall in $f$ may come from cash in hand. Why does that matter for judging the reform?
  \item What would you measure to tell the two channels apart?
  \item Which flow of the model would a firing cost act on instead?
  \item State the result: by how much does the reform raise the steady-state rate?
\end{enumerate}
\end{problem}
